% =============================================================================
\section{Related Work}
\label{sec:related}
% =============================================================================

\subsection{Single-Image Panorama Generation}

Extending a narrow-field photograph to a full $360^\circ$ environment is the
first step in most single-image scene systems. Diffusion models adapted to
panoramic geometry are the dominant approach.
CubeDiff~\cite{kalischek2025cubediffrepurposingdiffusionbasedimage} repurposes a
perspective latent diffusion model for panorama generation by operating on
cubemap faces with synchronized attention.
DreamCube~\cite{dreamcube,dreamcube_arxiv} extends this to RGB-D panoramas
through multi-plane synchronization, producing geometry alongside appearance.
PanoDreamer~\cite{paliwal2024panodreamer} takes an optimization-based route to
coherent single-image $360^\circ$ scenes, and
LayerPano3D~\cite{yang2024layerpano3d} decomposes the panorama into depth layers
to support larger viewpoint changes.
Our system supports several such backends, defaulting to a
FLUX-based~\cite{esser2024scalingrectifiedflowtransformers} outpainting path,
with WorldGen~\cite{worldgen2025ziyangxie} and DreamCube as alternatives. The
relevant distinction is that we treat the panorama as an intermediate
representation rather than as the deliverable.

\subsection{Single-Image 3D Scene Generation}

A growing body of work targets navigable 3D scenes from one image.
SphericalDreamer~\cite{schnepf2026sphericaldreamer} generates navigable immersive
worlds; SonoWorld~\cite{jin2026sonoworld} extends single-image scene generation to
audio-visual output; 360Anything~\cite{wu360anything} lifts images and videos
to $360^\circ$ geometry-free representations; and
WorldGen~\cite{worldgen2025ziyangxie} produces complete 3D scenes rapidly.
These systems generally optimize for visual fidelity from viewpoints near the
original camera. Into Frame instead prioritizes \emph{physical decomposition}: the
output is not one geometry proxy but a set of typed entities with distinct
behavior, a collidable terrain mesh, individually placed object meshes, an
instanced ground-cover population and an animated water surface. Each is required
before a scene can be dynamic rather than merely viewable.

\subsection{Depth Estimation and Panoramic Depth}

Monocular depth estimation underpins all geometry in the pipeline. We use
Depth Anything 3~\cite{lin2025depth3recoveringvisual} for perspective depth and
Depth Any Panoramas~\cite{lin2025dap} for equirectangular depth, the latter
avoiding the distortion artifacts that arise when a perspective model is applied
to equirectangular imagery. A recurring difficulty, which we address explicitly
in \cref{sec:panorama}, is that monocular depth is relative and saturates: real
distant terrain and sky receive indistinguishable predictions, so calibrated depth
at a ridge crest is an extrapolation rather than a measurement.

\subsection{Terrain Modeling and Heightmap Synthesis}

Recovering terrain from imagery connects to a long line of graphics work.
Pixels2Peaks~\cite{jain:hal-05617061} converts terrain images to heightmaps
directly; terrain modeling from feature primitives~\cite{genevaux:hal-01257198}
builds terrain from sparse user-specified features;
Earthbender~\cite{10.1145/3769047.3769053} offers interactive stylistic heightmap
generation; and WorldBrush~\cite{10.1145/2766975} synthesizes procedural worlds
from examples, including the distribution of scene elements. For physical
plausibility of unobserved relief we use
Landlab~\cite{barnhart2020landlab,hobley2017creative,Hutton_landlab_2020}, an
Earth-surface dynamics toolkit, applying fluvial incision so that invented terrain
obeys drainage structure rather than reading as smooth interpolation or undirected
noise.

\subsection{Texture Synthesis}

Projecting a panorama onto ground fails at grazing incidence: a single panorama
row covers meters of ground at distance. Pattern-based
texturing~\cite{neyret:inria-00537511} addresses exactly this class of problem by
assembling surfaces from patterns rather than resampling a projection, and our
terrain texturing follows that idea using patches extracted from the real
photograph. Where synthetic material is unavoidable we use
FLUX~\cite{esser2024scalingrectifiedflowtransformers} and latent diffusion
models~\cite{Rombach_2022_CVPR}, with
Swin2SR~\cite{conde2022swin2srswinv2transformercompressed} for super-resolution.

\subsection{Segmentation, Detection and Recognition}

Object decomposition uses SAM~2~\cite{ravi2024sam2segmentimages} for segmentation
and video tracking, Grounding DINO~\cite{liu2024groundingdino} for
open-vocabulary detection and label verification,
CLIP~\cite{radford2021learningtransferablevisual} for embedding-based
classification,
BLIP~\cite{li2022blipbootstrappinglanguageimagepretraining} and
Florence-2~\cite{xiao2024florence2} for captioning,
Recognize Anything~\cite{zhang2023recognize,huang2023tag2text} and its open-set
successor RAM++~\cite{huang2023open} for scene tagging,
SegFormer~\cite{DBLP:journals/corr/abs-2105-15203} for semantic region typing, and
BiRefNet~\cite{zheng2024birefnet} for high-resolution matting. Each solves a
similar but subtly different problem; we combine them rather than choosing between
them, using their agreement where they overlap to build a more reliable reading of
the scene than any one provides alone.

\subsection{Object Reconstruction and Inpainting}

Individual objects are reconstructed with
SAM~3D~\cite{sam3dteam2025sam3d3dfyimages}, with
SPAR3D~\cite{huang2025spar3dstablepointawarereconstruction} and
TRELLIS~\cite{xiang2025trellis2} as alternatives.
Removing foreground objects so the ground behind them can be reconstructed uses
inpainting: LaMa~\cite{suvorov2021resolution},
MAT~\cite{li2022mat}, Inpaint Anything~\cite{yu2023inpaint}, and
ObjectClear~\cite{zhao2026objectclear}, the last targeting removal of both objects
and their effects.

\subsection{Lighting}

Relighting reconstructed content so that it matches the original photograph
requires recovering the scene illumination. We use
LuxDiT~\cite{liang2025luxdit} for HDR lighting estimation
and IntrinsicDiffusion~\cite{Luo2024IntrinsicDiffusion} to separate albedo from
shading, so that assets carry material color rather than baked-in sunlight and
can be lit by our own lighting algorithm.
